\documentclass[aspectratio=1610,fontset=none]{ctexbeamer}
\usepackage[T1]{fontenc}

\title{Subliming 主题 Beamer 示例}
\author{杨曜堃\ \ 地球科学与信息物理学院}
\date{February 19, 2025}

% 添加封面上的大学校标，这里使用的是白色中南大学校标，并设置透明度为 30%
% 如果你想使用校徽，可以在 beamerinnerthemeSUBLIMING.sty 文件中修改，找到对应的注释即可
% 我的意思是，四个字的那种叫校徽，圆圆的那个叫校标
\titlegraphic{\includegraphics[width=8cm]{assets/csu-logo.pdf}}

\usepackage{styles/SUBLIMINGmacros}
\usefolder{styles}
\usetheme[style=csublue]{SUBLIMING} % 在这里更换颜色主题风格
% 支持 csu blue 和 nupurple 两种风格

\begin{document}

    \begin{frame}
        \titlepage
    \end{frame}

    \begin{frame}{介绍 Introduction}
        Subliming 主题是一个为中南大学设计的 Beamer 主题，也可以通过调整模板中的颜色、字体和背景以适应其他大学的风格。 
        \\ ~ \\
        这个主题是在 Northwestern University 的 Aaron Wolf 制作的 \href{https://github.com/aarondwolf/wildcatwildcat}{Wildcat} 主题的基础上修改设计的，更适合国内高校。同其一样，Subliming 主题也遵循 \href{http://www.gnu.org/licenses/}{GNU 通用公共许可证}，这意味着任何人都可以免费自由使用这个主题！
        \\ ~ \\
        你也可以自定义这个主题！我在各个包中增加了中文注释，以便你更好地理解这个主题的设计。如果有任何问题或建议，欢迎向我们的 \href{https://github.com/PourRevenir/Subliming}{GitHub 仓库}提交 \href{https://github.com/PourRevenir/Subliming/issues}{issue} 或 \href{https://github.com/PourRevenir/Subliming/pulls}{pull request}！
    \end{frame}

    \begin{frame}{目录 Contents}
        \tableofcontents
    \end{frame}

    \section{1 颜色和样式}

    \begin{frame}{彩色盒子示例 (tcolorbox)}
        这里提供了三种彩色盒子，你可以通过 \texttt{\textbackslash begin\{tblock\}} 环境创建. 你也可以使用 \texttt{talert} 和 \texttt{texample} 这两种环境分别生成其他颜色的盒子。
        \begin{tblock}{定理}
            这里是一个彩色盒子，颜色将和主题色一致。
        \end{tblock}
        \begin{talert}{说明}
            这里是一个说明，你也可以动手修改，比如制作一个可以自动编号的定理环境。
        \end{talert}
        \begin{texample}{例 1-1}
            这里是一个例子，其实只是一个绿色盒子。
        \end{texample}
    \end{frame}

    \begin{frame}[fragile]{可自定义颜色的盒子示例 (Facetbox)}
        使用\texttt{tfacetbox} 可以自定义盒子的颜色。 这只适用于非原色（红、绿、蓝），因为不容易进行阴影处理。
        \begin{tfacetbox}[nudarkyellow]{自定义颜色 Facet 盒子}
            这是一个 tfacetbox 盒子. 可通过以下代码创建:
            \begin{verbatim}
    \begin{tfacetbox}[nudarkyellow]{自定义颜色 Facet 盒子}
        这是一个 tfacetbox 盒子...
    \end{tfacetbox}
            \end{verbatim}
        \end{tfacetbox}
    \end{frame}

    \begin{frame}
        \frametitle{彩色盒子示例 (Default)}
        你也可以使用 Beamer 提供的默认文本盒子。默认样式没有圆角和阴影。
        \begin{block}{蓝色}
            这里是一个蓝色盒子
        \end{block}
        \begin{alertblock}{红色}
            这里是一个红色盒子
        \end{alertblock}
        \begin{exampleblock}{绿色}
            这里是一个绿色盒子
        \end{exampleblock}
        \end{frame}

    % \begin{frame}
    %     \frametitle{颜色 Colors}
    %     The theme has a few Northwestern-specific colors defined, which you can use in your slides. These are:
    %     \begin{columns}
    %     \column{0.3\textwidth}
    %     \begin{itemize}
    %         \item[$\textcolor{nupurple}{\bullet}$] \textcolor{nupurple}{nupurple} 
    %         \item[$\textcolor{nupurple90}{\bullet}$] \textcolor{nupurple90}{nupurple90} 
    %         \item[$\textcolor{nupurple80}{\bullet}$] \textcolor{nupurple80}{nupurple80}
    %         \item[$\textcolor{nupurple70}{\bullet}$] \textcolor{nupurple70}{nupurple70}
    %         \item[$\textcolor{nupurple60}{\bullet}$] \textcolor{nupurple60}{nupurple60}
    %         \item[$\textcolor{nupurple50}{\bullet}$] \textcolor{nupurple50}{nupurple50}
    %         \item[$\textcolor{nupurple40}{\bullet}$] \textcolor{nupurple40}{nupurple40}
    %         \item[$\textcolor{nupurple30}{\bullet}$] \textcolor{nupurple30}{nupurple30}
    %         \item[$\textcolor{nupurple20}{\bullet}$] \textcolor{nupurple20}{nupurple20}
    %         \item[$\textcolor{nupurple10}{\bullet}$] \textcolor{nupurple10}{nupurple10}
    %     \end{itemize}
    %     \column{0.3\textwidth}
    %     \begin{itemize}
    %         \item[$\textcolor{nupurple160}{\bullet}$] \textcolor{nupurple160}{nupurple160}
    %         \item[$\textcolor{nupurple150}{\bullet}$] \textcolor{nupurple150}{nupurple150}
    %         \item[$\textcolor{nupurple140}{\bullet}$] \textcolor{nupurple140}{nupurple140}
    %         \item[$\textcolor{nupurple130}{\bullet}$] \textcolor{nupurple130}{nupurple130}
    %         \item[$\textcolor{nupurple120}{\bullet}$] \textcolor{nupurple120}{nupurple120}
    %         \item[$\textcolor{nupurple110}{\bullet}$] \textcolor{nupurple110}{nupurple110}
    %         \item[$\textcolor{nurichblack}{\bullet}$] \textcolor{nurichblack}{nurichblack}
    %         \item[$\textcolor{nubrightgreen}{\bullet}$] \textcolor{nubrightgreen}{nubrightgreen}
    %         \item[$\textcolor{nubrightteal}{\bullet}$] \textcolor{nubrightteal}{nubrightteal}
    %         \item[$\textcolor{nubrightblue}{\bullet}$] \textcolor{nubrightblue}{nubrightblue}
    %     \end{itemize}
    %     \column{0.3\textwidth}
    %     \begin{itemize}
    %         \item[$\textcolor{nubrightyellow}{\bullet}$] \textcolor{nubrightyellow}{nubrightyellow}
    %         \item[$\textcolor{nubrightorange}{\bullet}$] \textcolor{nubrightorange}{nubrightorange}
    %         \item[$\textcolor{nubrightred}{\bullet}$] \textcolor{nubrightred}{nubrightred}
    %         \item[$\textcolor{nudarkgreen}{\bullet}$] \textcolor{nudarkgreen}{nudarkgreen}
    %         \item[$\textcolor{nudarkteal}{\bullet}$] \textcolor{nudarkteal}{nudarkteal}
    %         \item[$\textcolor{nudarkblue}{\bullet}$] \textcolor{nudarkblue}{nudarkblue}
    %         \item[$\textcolor{nudarkyellow}{\bullet}$] \textcolor{nudarkyellow}{nudarkyellow}
    %         \item[$\textcolor{nudarkorange}{\bullet}$] \textcolor{nudarkorange}{nudarkorange}
    %         \item[$\textcolor{nudarkred}{\bullet}$] \textcolor{nudarkred}{nudarkred}
    %     \end{itemize}
    %     \end{columns}
    % \end{frame}

    % \begin{frame}{Modifying Main Colors (I)}
    %     You can change the main colors of the theme by redefining the following colors in your preamble. You can use any color you want, but the theme is designed to work best with shades of the Northwestern purple.
    %     \begin{itemize}
    %         \item \texttt{wcprimary} (main color)
    %         \item \texttt{wcprimary10} (main color, 10\% shade) through \texttt{wcprimary40} (main color, 40\% shade)
    %         \item \texttt{wcprimary110} (main color, 110\% shade) through \texttt{wcprimary160} (main color, 160\% shade)
    %         \item \texttt{wcalerted} (alert color)
    %         \item \texttt{wcexample} (example color)
    %     \end{itemize}
    % The shades of wcprimary are used for the background of facets. There are wcprimary10, wcprimary20, ..., wcprimary160, but only 10-40 and 110-140 are used.
    % \end{frame}

    % \begin{frame}[fragile]
    %     \frametitle{Modifying Main Colors (II)}
    %     For example, to modify the main color to be a shade of blue, you could use the following code in your preamble (after loading the theme):
    %     \begin{verbatim}
    %         \definecolor{wcprimary}{RGB}{0,53,107}
    %         \definecolor{wcprimary140}{RGB}{0, 34, 70}
    %         \definecolor{wcprimary130}{RGB}{0, 40, 80}
    %         \definecolor{wcprimary120}{RGB}{0, 45, 91}
    %         \definecolor{wcprimary110}{RGB}{0, 50, 102}
    %         \definecolor{wcprimary40}{RGB}{153, 174, 196}
    %         \definecolor{wcprimary30}{RGB}{179, 194, 211}
    %         \definecolor{wcprimary20}{RGB}{204, 215, 225}
    %         \definecolor{wcprimary10}{RGB}{230, 235, 240}
    %     \end{verbatim}
    %     This sets the main color to Yale Blue, and all the shades used in the facet pattern. The preambles of the example slides in this document show how to modify the colors.
    % \end{frame}

    % \begin{frame}[fragile]
    %     \frametitle{Using Colors in Graphs}
    %     %You can use the colors in your graphs in Python, R, and Stata using the following commands:
    %     \begin{tblock}{Python Code}
    %         \footnotesize
    %         \begin{verbatim}
    % import matplotlib.pyplot as plt
    % import pandas as pd
    % plt.style.use('ggplot')
    % # Create colors
    % wcprimary = (78/255,42/255,132/255)
    % df = pd.read_stata("source/graphs/auto.dta")
    % plt.scatter(df["weight"], df["mpg"], color=wcprimary)
    % plt.xlabel("Weight")
    % plt.ylabel("MPG")
    % plt.savefig("source/graphs/plot-python.pdf")
    % plt.show()
    %         \end{verbatim}
    %     \end{tblock}
    % \end{frame}

    % \begin{frame}[fragile]
    %     \frametitle{Using Colors in Graphs}
    %     \begin{tblock}{Stata Code}
    %         \footnotesize
    %         \begin{verbatim}
    % *ssc install schemepack, replace
    % sysuse auto, clear
    % local wcprimary "78 42 132"

    % twoway (scatter mpg weight, mcolor("`wcprimary'")), ///
    %         scheme(gg_tableau)
            
    % graph export "plot-stata.pdf", replace
    %         \end{verbatim}
    %     \end{tblock}
    % \end{frame}

    % \begin{frame}{Font Style: Default}
    %     The theme uses the official Northwestern fonts, which are Campton (for titles) and Akkurat Pro (for copy). These are not free fonts, so you will need to purchase (and install) them if you want to use them.
    %     \\ ~ \\
    %     The wildcat theme uses the OTF files directly, so as long as you have the OTFs in a similar relative directory as the demo files, it should work (this includes on Overleaf). If you have installed them, you can look at the \texttt{\textbackslash beamerfontthemewildcat-installed.sty} file to see how to use the installed fonts.
    %     \\ ~ \\
    %     Note: You need to us XeLaTeX or LuaLaTeX to compile in order to use custom fonts. If you use PDFLaTeX, you will get an error.
    % \end{frame}

    % \begin{frame}{Font Style: Overleaf \& Local}
    %     There is also a font theme called \texttt{wildcat-overleaf} which will load fonts available from Overleaf (as well as local \LaTeX ~ installations). To use this, you just need to specify in your preamble the following command \textit{after} loading the theme: 
    %     \texttt{\textbackslash usefonttheme\{wildcat-overleaf\}}
    % \end{frame}

    % \begin{frame}[fragile]{Font Style: Default in Overleaf}
    %     An alternative is to upload the OTF or TTF files to your Overleaf directory, and then change the \texttt{beamerfontthemewildcat.sty} files to specify those font files rather than the font family generally. Alternatively, simply add the lines of code directly to your preamble \textit{after} loading the theme:
    %     \footnotesize
    %     \begin{verbatim}
    % % Create new font families from OTF fonts
    % \newfontfamily\CamptonMedium[Path=fonts/Campton/, Extension=.otf]{Campton Medium}
    % \newfontfamily\CamptonLight[Path=fonts/Campton/, Extension=.otf]{Campton Light}

    % % Change title font to Campton
    % \setbeamerfont{title}{family=\CamptonMedium, size=\LARGE}
    % \setbeamerfont{author}{family=\CamptonLight, size=\small}
    % \setbeamerfont{frametitle}{family=\CamptonMedium, size=\Large}
    % \setbeamerfont{section}{family=\CamptonLight, size=\LARGE}
    %     \end{verbatim}
    % \end{frame}

    % \begin{frame}{Font Style: Examples}
    %     Here is what the body font will look like under normal usage:
    %     \begin{itemize}
    %         \item Regular
    %         \item \textit{Italic}
    %         \item \textbf{Bold}
    %         \item \textbf{\textit{Bold Italic}}
    %         \item \alert{Alert}
    %         \item \alert{\textit{Alert Italic}}
    %         \item Math: $$e = \lim_{n \rightarrow \infty}\left(1+\frac{1}{n}\right)^{n}$$
    %     \end{itemize}
    % \end{frame}

    % \begin{frame}{Sections}
    %     You can use \texttt{\textbackslash section\{\}} to group slides into sections. This will automatically add a section title slide at the beginning of each section.
    %     \\ ~ \\
    %     You can turn this off by adding \texttt{\textbackslash AtBeginSection\{\}} to your preamble (this tells beamer to do precisely nothing when a new section starts).
    % \end{frame}


    

    % \begin{frame}[fragile]{Advanced: Altering the background pattern}
    %     The background pattern is created using the \texttt{tikz} package. You can modify the background pattern by changing the \texttt{beamerinnerthemewildcat.sty} file. The background pattern is created in the \texttt{\textbackslash bgpattern} command defined at the top of the file.
    %     \\ ~ \\
    %     You don't necessarily need to change the file itself. You can just re-define those commands in your preamble. For example, if you wanted to change the background pattern to be just the main color, you could use the following code in your preamble: \\
        
    %     \scriptsize
    %     \begin{verbatim}
    % \renewcommand{\bgpattern}{
    %     \draw[color=wcprimary,fill=wcprimary] (0,0) rectangle (\paperwidth,\paperheight);
    % }
    %     \end{verbatim} 
    % \end{frame}




    % % Blank Frame
    % \begin{frame}[plain]{}
    %     Here is a blank frame. You can use this to add a blank page to your presentation. You can specify this via:
    %     \\ ~ \\
    %     \texttt{\textbackslash begin\{frame\}[plain]\{\}} \\
    %     \texttt{...} \\
    %     \texttt{\textbackslash end\{frame\}}
    % \end{frame}

    % % Standout frame
    % \begin{frame}{Standout Slides}
    %     You can use the \texttt{standout} command to create a slide with a large font and no title. This is useful for a slide that you want to stand out, but which doesn't need a title.
    %     \\ ~ \\
    %     \texttt{\textbackslash standout\{Questions?\}}

    % \end{frame}

    \standout{谢谢 Thank You}

\end{document}